\begin{table*}[ht!]
\center
\small
\begin{tabular}{l  c  c  l  l l}
\toprule 
% \small
% \small
\textbf{Task}  & $|\boldsymbol{q}|$ &  $|\mathcal{C}|$ & \textbf{Domain} & \textbf{Query} & \textbf{Gold documents} \\  
\midrule
Ambig QA~\cite{min-etal-2020-ambigqa} & 1,172 & 18,809 & Wikipedia & question & duplicated question  \\
WIKIQA~\cite{yang-etal-2015-wikiqa} & 369 & 26,196 & Wikipedia & question & answer sentence \\
\midrule
SciFact~\cite{wadden-etal-2020-fact} & 300 & 5183 & Science & claim & scientific paper paragraph \\ 
\midrule
GooAQ-Technical~\cite{khashabi-etal-2021-gooaq-open} & 1,000 & 4,086 & Technical & question & StackOverflow answer \\
LinkSo-Python~\cite{linkso} & 1,000 & 485,413 &  Technical & question & StackOverflow question \\
CodeSearchNet-Python~\cite{husain2019codesearchnet} & 1,000 & 457,414 & Code & comment & Python code \\
% CodeSearch-java & 1k & 500k  & Code & comment & Java code \\
\bottomrule
\end{tabular}
\caption{The \task evaluation. Example pairs of queries and documents are shown in Table~\ref{tab:cross_task_eval_examples}. In addition to the corpora listed above, we add the Natural Questions corpus data from BEIR~\cite{thakur2021beir}, which consists of 2,681,468 documents from Wikipedia. 
}
\label{tab:cross_task_eval}
\end{table*}

\section{Experiments}
{
We evaluate
% our instruction-following retrieval systems 
\model on zero-shot retrieval (Section~\ref{sec:beir_lotte}) and in a challenging new evaluation setup, \task (Cross-task Cross-domain Retrieval; detailed in Section~\ref{sec:evaluatiom_benchmark}), comparing them with state-of-the-art models described in Section~\ref{sec:baselines}. 
}


\subsection{Zero-shot Retrieval Evaluations}
\label{sec:beir_lotte}
To evaluate the models' ability to perform zero-shot transfer via {\it instructions}, we run experiments on two widely used zero-shot transfer evaluation benchmarks: {\bf BEIR}~\cite{thakur2021beir} and {\bf LOTTE-pooled}~\cite{santhanam-etal-2022-colbertv2}.
{
Notably, none of the evaluation datasets and instructions overlap with \rib. 
Moreover, many tasks differ significantly from tasks used during training (e.g., argument retrieval)}.

{\bf BEIR} is a collection of diverse retrieval tasks in multiple domains (e.g., Wikipedia, biomedical) where the retrieval target is restricted to the target corpus in a single domain. We used publicly available datasets. Following \citet{promptagator}, we exclude Natural Questions, MS MARCO, HotpotQA, FEVER and CQADupStack from our evaluation targets for fair comparison since they are included either in encoders' pretraining or in \rib. 

{\bf LOTTE-Search} samples GooAQ~\cite{khashabi-etal-2021-gooaq-open} questions whose answers come from certain forums in StackExchange. 
We evaluate our model in the pooled setup, where documents come from forums in diverse domains (e.g., cooking, technical). GooAQ is not included in our training set. 
In LOTTE, our instructions specify which {\it forum domains} our system should retrieve evidence from (e.g., ``Retrieve a {\it cooking} StackExchange forum post that answers this question''). 

\paragraph{Metrics.}
Following \citet{thakur2021beir}, for BEIR, we use NDCG@10 as our primary metric on BEIR. For LOTTE-pooled, we use Success@5 ($=$ Recall@5) as our primary metric, as in the original paper~\cite{santhanam-etal-2022-colbertv2}.

\subsection{\task: Cross-task Cross-domain Retrieval Evaluation} 
\label{sec:evaluatiom_benchmark}
{
{
Normal retrieval benchmarks often assume that a system must deal with only a single intent and a closed corpus, which may oversimplify real-world scenarios: users' intents can be diverse, requiring searching in a truly open-domain environment that includes diverse documents~\cite{piktus2021web}. 
}
We introduce a more realistic evaluation setup, \task, where several retrieval tasks with different intents are pooled to form a single retrieval target containing diverse documents. This setup requires a system to adapt to a new task in a zero-shot manner and to model users' intents expressed in natural languages, to find documents aligning their expectations from an open-domain corpora. 
}

\paragraph{Tasks and queries.}
Our \task evaluation covers six datasets across three domains, namely, Wikipedia, Science, and Technical (Table~\ref{tab:cross_task_eval}) domains. 
The key challenge here includes datasets with different search intents that may not always be obvious from the queries alone. 
% Even when queries look similar the expected retrieval results are different given different implicit intents. 

\paragraph{A pooled corpus.}
For the primary \emph{pooled} setup, we combine all documents from different tasks and the BEIR NQ Wikipedia corpus to form a single retrieval corpus, consisting of approximately 3.7 million documents. 
We also report the simplified \emph{closed} setup performance as an oracle setup, where a system retrieves only from the original task corpus, which is similar to BEIR. 

\paragraph{{Metrics}.}
We report NDCG@10 on both pooled and closed setups for each task. In addition, we evaluate the performance gap between the closed and pooled setups and refer to it as \emph{robustness}.
A smaller gap means that the model is distracted less by the documents from undesirable corpora.


\begin{table*}[t!]
\centering
\footnotesize
\addtolength{\tabcolsep}{-4.25pt}  
\begin{tabular}{@{}l|a|a | a| cccccccccb|c@{}}\toprule
 & \multicolumn{3}{c|}{model size \& rerank}   &\multicolumn{10}{c}{\textbf{BEIR}} & \textbf{LOTTE} \\ \midrule
& Ret. & Gen. & $K$ &TREC & NFC & FQA & ARG& TOU & DBP & SCD & CLI&SCF & Avg. & Search-Pooled \\\midrule
BM 25 & 0 & 0 & 0 & 65.6  & 32.5  & 23.6 & 31.5 & 36.7 & 31.3 & 15.8 &  21.3&  66.5 & 36.0 & 48.3  \\
Contriever & 110M & 0 & 0 & 27.4 & 31.7 & 24.5 & 37.9 & 19.3 & 29.2 & 14.9 & 15.5 & 64.9 & 29.3 & 55.5 \\ 
UPR$^\dagger$ & 3B & 0 & 0 & 60.4 & 33.3 & 45.0 & 50.3 &  21.3 & 33.8 & 17.3 & 9.5 & 69.6 & 37.8  & --  \\ 
% SGPT-CE & 6.1 B & 0 & &  79.1 & 34.7 & 40.1 & 28.6 & 23.4 & 37.0 & 19.6 & 16.1 & 68.2 & 38.6 \\ 
\hline
Contriever (MS) & 110M & 0 & 0 & 59.6 &	32.8 &	32.9 &	44.6 &	23.0 & 41.3	& 16.5 & 	23.7 & 	67.7 & 	38.0 &  66.0  \\
Contriever+CE$^\dagger$& 133M & 0&  100&  70.1 &34.4 & 	36.7 &	41.3&	{ 29.8} &	47.1 &	17.1 &	25.8 &	69.2 &	41.3  & 73.5  \\
ColBERT-v2 & 110M &  0& 0 &  73.8 & 33.8 & 35.6 & 47.9 & 26.3 & 44.6 & 15.8 & 17.6 & 69.3 & 40.5 &   71.6 \\
BM25 + MonoT5 (3B)$^\dagger$ & 3B & 0& 1000& {\bf 79.6} & {\bf 38.4} & {\bf 51.2} &	28.8& 20.0 & {\bf 47.8} &	18.4 &	28.9 & 	{\bf 77.7} & 	43.4  & --   \\
GTR-11B& 4.8B & 0& 0 &  50.1 & 34.2 & 46.7 & {\bf 54.0} & 25.6 & 40.8& 16.1 & 27.0 &  63.5 & 39.8 & --  \\
SGPT-6.8B & 6.8B & 0 & 0 & 87.3 & 36.2 & 37.2 & 51.4 & 25.4 & 39.9 & {\bf 19.7} & 30.5 & 74.7 & 44.7 & --   \\
\midrule
GPL &  66M$\times$9 & 220M & 0 & 72.6 & -- & 32.8 & -- & -- &  -- & -- & -- & 66.4 & -- & --\\
Promptagator & 110M $\times$9 & 175B & 0 & 72.7 & 33.4 & 40.4 & 53.8 & 26.6 & 36.4 & 16.3 & 21.4 & 62.3 & 40.4 & --  \\
Promptagator (rank)$^\dagger$ & 220M $\times$9 & 175B & 200 & 76.0 &	36.0 & 	45.9 & 	53.1 & 	27.8  & 	41.3 & 	{ 19.1} &	22.6 & 	73.6 & 	43.9
& -- \\
\midrule
{\bf \model-dual} &110M &0 & 0 & 62.6 & 33.7 & 33.7 & 48.9 & 20.1 & 41.5 & 14.2 & 13.8 & 69.0 & 37.4 & 56.8  \\
{\bf \model-full (T0-3B)}$^\dagger$ &1.5B &0 & 100 &  71.7 & 34.0 & 42.2& 	49.8 & 	{\bf 31.2} & 	45.1 & 	17.5 & 	{ 30.0} & 	75.8 & 	{ 44.1}	& {\bf 75.7}  \\ 
{\bf \model-full (FLAN-T5)}$^\dagger$ &1.5B &0 & 100 & { 72.8} & 33.4 & 41.8 & { 51.5} & 24.9 & {46.8} & { 18.7} & {\bf  35.4} & {\bf 77.7} & {\bf 44.8}	& 73.1 \\
\bottomrule
\end{tabular}
    \caption{Zero-shot retrieval results on BEIR and LOTTE-Search (pooled). $\dagger$ indicates the models using cross-encoder-based reranking models. The first group of models use no labeled data during training. The second group uses MS MARCO at training time but have no customized task-specific data. The third group trains individual retrieval systems using automatically generated data. TREC, NFC, FQA, ARG, TOU, DBP, SCD, CLI, SCF indicates TREC-COVID~\cite{10.1145/3451964.3451965}, FIQA~\cite{10.1145/3184558.3192301}, NF Corpus~\cite{nfcorpus}, Arguana~\cite{wachsmuth-etal-2018-retrieval}, Touche-2020~\cite{touch}, DBPedia~\cite{dbpedia}, SciDocs~\cite{cohan-etal-2020-specter}, Climate- Fever~\cite{climate_fever}, and SciFact~\cite{wadden-etal-2020-fact}, respectively. ``$\times 9$'' of GPL, Promptagator means that those models train customized models for each of the datasets.  
    }
    \label{tab:beir}
\end{table*}
\begin{table*}[t!]
\centering
\footnotesize
\addtolength{\tabcolsep}{-2.pt}  
\begin{tabular}{l|cc|cc|cc|cc|cc|cc|bb| b}\toprule
& \multicolumn{2}{c}{AMB}  &
\multicolumn{2}{c}{WQA} &  \multicolumn{2}{c}{SCF}  &  \multicolumn{2}{c}{GAT} & \multicolumn{2}{c}{LSO} & \multicolumn{2}{c}{CSP} & \multicolumn{2}{b}{Avg.} &  $\Delta$ \\
&cl&pl&cl&pl&cl&pl&cl&pl&cl&pl&cl&pl&cl&pl & cl$-$pl\\
\midrule
% BM25 &  & & &&  & & & & & & &   \\
Contriever & {\bf 96.8} & { 93.8}  & 80.9& 54.1 &  67.7 & 57.4  & 73.2 & 59.8 & { 28.0} & 26.7 & 36.7 &  36.1 & 63.9 & 54.6 & 9.3 \\
% Contr & \\\hline
Contriever+CE & 96.6  & 47.4  & 78.2 & 58.4 & 69.1 & 61.7 &  75.4 &66.0  & {\bf 32.1} & {\bf 31.4} & 42.0 & 40.2  & 65.5 & 50.9 & 13.4 \\
\model-dual &  96.3 &  {\bf 95.3}  &  80.2 &   {\bf 63.1} & 70.1  & {66.2 } & 75.0 & 65.0 & 23.0 & 23.4 & 31.3 & 31.3 & 60.5  &  53.6 & {\bf 6.9} \\
{\bf \model-full (T0)} & 91.1 &  90.5 & { 82.1} &  52.5  &  { 74.7} & {66.2}  & {\bf 80.5} & {\bf 68.6} & 25.1 & 24.9 & { 51.4}  & {\bf 51.4} & { 67.5} & {\bf 59.1} & 8.4   \\
{\bf \model-full (FLAN)} & 94.0 &  89.6 &  8\bf 6.9 & 55.9 &  \bf 77.4 & \bf 66.3 & 78.3 & 66.7  & 18.1 & 18.4 & \bf 51.8 & 50.1 & {\bf 67.7} & 57.8 &  9.9 \\
\bottomrule
\end{tabular}
    \caption{Cross-task retrieval results.  $\Delta$ shows the performance gap between the averaged pooled performance and the averaged closed performance. AMB, WQA, SCF, GAT, LSO, CSP denote AmbigQA, WikiQA, SciFact, GooAQ-Technical, LinkSO-Python, and CodeSearchNet-Python, respectively. 
    }
    \label{tab:cross_task_results}
\end{table*}
\subsection{Baselines}
\label{sec:baselines}
We compare \model with various state-of-the-art methods. 
The first group we consider are unsupervised models that are not trained or trained only on unlabeled text;  
these include {\bf Contriever}~\cite{izacard2022unsupervised} and {\bf BM25}. 
We also compare \model with {\bf UPR}~\cite{sachan2022improving}, which reranks the Contriever results using a pretrained T0-3B. 

The second group trains retrievers and rerankers on MS MARCO or a few large-scale datasets and directly transfers them to new tasks with no adaptations, including {\bf MonoT5}~\cite{nogueira-etal-2020-document}, {\bf Contriever-MS MARCO} and {\bf Contriever-MS MARCO + Cross Encoder (CE)}, {\bf ColBERT v2}~\cite{santhanam-etal-2022-colbertv2}, {\bf GTR}~\cite{ni2021large} and {\bf SGPT-BE}~\cite{muennighoff2022sgpt}. 

The final group of models is specialized retrievers trained for each task on additional task-specific data that was automatically generated given the target corpus. 
{\bf Promptagator}~\cite{promptagator} generates large amount of in-domain data using FLAN~\cite{wei2022finetuned}, and {\bf GPL}~\cite{wang-etal-2022-gpl} generates them using DocT5Query~\cite{doc2query}. 

\subsection{Experimental Settings}
We initialize \model-full from the T0-3B encoder~\cite{sanh2022multitask} {and FLAN-T5 encoder~\cite{flan_palm}}. We sample positive and negative passages with a 1:4 ratio. 
We train \model-full up to 10k steps and take the best checkpoint based on development split performance. 
We initialize \model-dual from a Contriever-MS MARCO checkpoint~\cite{izacard2022unsupervised} and train up 30k steps. 
Per-GPU batch size is 16, and for each positive document, we sample in total 5 negative passages; 90\% of them are randomly sampled from $\mathcal{D}$, and 10\% are sampled from $\boldsymbol{d}^{\rm HD}$ and $\boldsymbol{d}^{\rm UF}$ in addition to the in-batch negative documents. 
We use 8 GPUs to train \model-full and 64 GPUs to train \model-dual.
To retrieve the initial document candidates for \model-full, we use Contriever-MS MARCO and rerank the top 100 documents.\footnote{We found that combining \model-full with the original Contreiver performs better than combining \model-full with \model-dual, possibly because \model-full uses the hard negative samples retrieved by Contriever's top-retrieved results. }
{Table~\ref{tab:list_of_instructions_beir} shows the full list of instructions for evaluations}. 
More details are in Appendix~\ref{sec:hyperparameters}. 
